% GENERATED FILE. Edit canonical lessons, then run npm run generate:books.
% canonical-source-hash: fnv1a64:1550bc5decd18038
% canonical-lessons: HI-C131-tauliya, HI-C131-jharu, HI-C131-dibba, HI-C131-thaila, HI-C131-mombatti

\chapter{Towel, Broom, Box, Bag, Candle}
\label{ch:hi-towel-broom-box-bag-candle}
\hlchaptermodality{131}

\begin{chapteropening}
\textbf{By the end of this chapter:} \emph{I can say a towel, a broom, a box, a bag and a candle.}

Complete the last lesson of chapter 131: I can say a towel, a broom, a box, a bag and a candle.
\end{chapteropening}

\section[\texorpdfstring{\hi{तौलिया} (tauliyā)}{तौलिया (tauliyā)}]{\hi{तौलिया} (tauliyā) — a towel}

\label{lesson:HI-C131-tauliya}



\begin{quote}
Before the new one: say the Hindi for thread, then the Hindi for a comb.
\end{quote}

\subsection*{You'll want to know: \hi{तौलिया}}
\textbf{\hi{तौलिया}} — \emph{tauliyā} — \textquotedblleft{}a towel\textquotedblright{}.

\begin{cousinweb}
\textbf{\hi{तौलिया}} (\emph{tauliyā}) is usually traced to Portuguese \emph{toalha}, \textquotedblleft{}towel\textquotedblright{}.
\end{cousinweb}

\begin{grammarlens}[title={What you've built}]
One more word for this chapter.
\end{grammarlens}

\subsection*{Your turn}
\begin{itemize}
\raggedright
  \item \emph{Say it:} \emph{tauliyā}
  \item \emph{Say it:} \emph{tauliyā}, once more
  \item \emph{Say it:} say \emph{tauliyā}
  \item \emph{Recall:} say the Hindi for thread, then the Hindi for a comb, then say \emph{tauliyā} again
\end{itemize}

\begin{tcolorbox}[breakable,colback=teal!4,colframe=teal!35!black,title={Before you move on}]
What does \hi{तौलिया} mean? (\textquotedblleft{}A towel\textquotedblright{}.) Say it once more.
\end{tcolorbox}
\section[\texorpdfstring{\hi{झाड़ू} (jhāṛū)}{झाड़ू (jhāṛū)}]{\hi{झाड़ू} (jhāṛū) — a broom}

\label{lesson:HI-C131-jharu}



\begin{quote}
Before the new one: say the Hindi for a comb, then the Hindi for a towel.
\end{quote}

\subsection*{You'll want to know: \hi{झाड़ू}}
\textbf{\hi{झाड़ू}} — \emph{jhāṛū} — \textquotedblleft{}a broom\textquotedblright{}.

\begin{cousinweb}
\textbf{\hi{झाड़ू}} (\emph{jhāṛū}) is built on \emph{jhāṛnā}, \textquotedblleft{}to sweep, shake off\textquotedblright{}, which CDIAL lists under a reconstructed \emph{*jhāṭayati}, \textquotedblleft{}shakes down\textquotedblright{} (CDIAL 5328).
\end{cousinweb}

\begin{grammarlens}[title={What you've built}]
One more word for this chapter.
\end{grammarlens}

\subsection*{Your turn}
\begin{itemize}
\raggedright
  \item \emph{Say it:} \emph{jhāṛū}
  \item \emph{Say it:} \emph{jhāṛū}, once more
  \item \emph{Say it:} say \emph{jhāṛū}
  \item \emph{Recall:} say the Hindi for a comb, then the Hindi for a towel, then say \emph{jhāṛū} again
\end{itemize}

\begin{tcolorbox}[breakable,colback=teal!4,colframe=teal!35!black,title={Before you move on}]
What does \hi{झाड़ू} mean? (\textquotedblleft{}A broom\textquotedblright{}.) Say it once more.
\end{tcolorbox}
\section[\texorpdfstring{\hi{डिब्बा} (ḍibbā)}{डिब्बा (ḍibbā)}]{\hi{डिब्बा} (ḍibbā) — a box}

\label{lesson:HI-C131-dibba}



\begin{quote}
Before the new one: say the Hindi for a towel, then the Hindi for a broom.

Say \textbf{\hi{तीन}}, \textbf{\hi{चार}} and \textbf{\hi{पाँच}}, then the path behind the first: \emph{trīṇi}, then \emph{tiṇṇi}, then \textbf{\hi{तीन}}.
\end{quote}

\subsection*{You'll want to know: \hi{डिब्बा}}
\textbf{\hi{डिब्बा}} — \emph{ḍibbā} — \textquotedblleft{}a box\textquotedblright{}.

\begin{cousinweb}
\textbf{\hi{डिब्बा}} (\emph{ḍibbā}) belongs with Hindi \emph{ḍabbā} and \emph{ḍibbī}, \textquotedblleft{}box\textquotedblright{}, which CDIAL lists under a reconstructed \emph{*ḍabba}, \textquotedblleft{}box\textquotedblright{}, and its variant \emph{*ḍibba} (CDIAL 5528). CDIAL takes the group to be non-Aryan, perhaps Dravidian.
\end{cousinweb}

\begin{grammarlens}[title={What you've built}]
One more word for this chapter.
\end{grammarlens}

\subsection*{Your turn}
\begin{itemize}
\raggedright
  \item \emph{Say it:} \emph{ḍibbā}
  \item \emph{Say it:} \emph{ḍibbā}, once more
  \item \emph{Say it:} say \emph{ḍibbā}
  \item \emph{Recall:} say the Hindi for a towel, then the Hindi for a broom, then say \emph{ḍibbā} again
\end{itemize}

\begin{tcolorbox}[breakable,colback=teal!4,colframe=teal!35!black,title={Before you move on}]
What does \hi{डिब्बा} mean? (\textquotedblleft{}A box\textquotedblright{}.) Say it once more.
\end{tcolorbox}
\section[\texorpdfstring{\hi{थैला} (thailā)}{थैला (thailā)}]{\hi{थैला} (thailā) — a bag}

\label{lesson:HI-C131-thaila}



\begin{quote}
Before the new one: say the Hindi for a broom, then the Hindi for a box.
\end{quote}

\subsection*{You'll want to know: \hi{थैला}}
\textbf{\hi{थैला}} — \emph{thailā} — \textquotedblleft{}a bag\textquotedblright{}.

\begin{cousinweb}
\textbf{\hi{थैला}} (\emph{thailā}) is a tadbhava, a word inherited from Sanskrit and reshaped by sound change. It comes from \emph{sthavi}, \textquotedblleft{}sack, bag\textquotedblright{} (CDIAL 13746). The \emph{-l-} is an added suffix; the Sanskrit word is known only from the old word-lists.
\end{cousinweb}

\begin{grammarlens}[title={What you've built}]
One more word for this chapter.
\end{grammarlens}

\subsection*{Your turn}
\begin{itemize}
\raggedright
  \item \emph{Say it:} \emph{thailā}
  \item \emph{Say it:} \emph{thailā}, once more
  \item \emph{Say it:} say \emph{thailā}
  \item \emph{Recall:} say the Hindi for a broom, then the Hindi for a box, then say \emph{thailā} again
\end{itemize}

\begin{tcolorbox}[breakable,colback=teal!4,colframe=teal!35!black,title={Before you move on}]
What does \hi{थैला} mean? (\textquotedblleft{}A bag\textquotedblright{}.) Say it once more.
\end{tcolorbox}
\section[\texorpdfstring{\hi{मोमबत्ती} (mombattī)}{मोमबत्ती (mombattī)}]{\hi{मोमबत्ती} (mombattī) — a candle}

\label{lesson:HI-C131-mombatti}



\begin{quote}
Before the new one: say the Hindi for a box, then the Hindi for a bag.
\end{quote}

\subsection*{You'll want to know: \hi{मोमबत्ती}}
\textbf{\hi{मोमबत्ती}} — \emph{mombattī} — \textquotedblleft{}a candle\textquotedblright{}.

\begin{cousinweb}
\textbf{\hi{मोमबत्ती}} (\emph{mombattī}) is \textquotedblleft{}wax-wick\textquotedblright{}. \emph{mom}, \textquotedblleft{}wax\textquotedblright{}, is usually traced to Persian \emph{mom}, and \emph{battī}, \textquotedblleft{}wick\textquotedblright{}, comes from Sanskrit \emph{varti}, \textquotedblleft{}wick\textquotedblright{} (CDIAL 11359).
\end{cousinweb}

\begin{grammarlens}[title={What you've built}]
One more word for this chapter.
\end{grammarlens}

\subsection*{Your turn}
\begin{itemize}
\raggedright
  \item \emph{Say it:} \emph{mombattī}
  \item \emph{Say it:} \emph{mombattī}, once more
  \item \emph{Say it:} say \emph{mombattī}
  \item \emph{Recall:} say the Hindi for a box, then the Hindi for a bag, then say \emph{mombattī} again
\end{itemize}

\begin{tcolorbox}[breakable,colback=teal!4,colframe=teal!35!black,title={Before you move on}]
What does \hi{मोमबत्ती} mean? (\textquotedblleft{}A candle\textquotedblright{}.) Say it once more.
\end{tcolorbox}
